% Pollis New Methods: volatility (periodic GARCH) — draft 1 (11 October 2026)
% Compile with: latexmk -pdf nm_volatility.tex && latexmk -c
% Shared look for the Pollis New Methods documents (same fonts and colours as the brochure and the workflow handouts).
\documentclass[10pt,a4paper]{article}

\usepackage[a4paper,left=18mm,right=18mm,top=15mm,bottom=17mm,footskip=8mm]{geometry}
\usepackage[T1]{fontenc}
\usepackage[default,semibold]{sourcesanspro}
\usepackage{textcomp}
\usepackage[italic]{mathastext}
\usepackage{amsmath}
\usepackage{microtype}
\usepackage[table]{xcolor}
\usepackage{tikz}
\usepackage[most]{tcolorbox}
\usepackage{fontawesome5}
\usepackage{tabularx,booktabs}
\usepackage{enumitem}
\usepackage{fancyhdr}
\usepackage[hidelinks]{hyperref}

% ---------------------------------------------------------------- colours (as in the brochure)
\definecolor{navy}{HTML}{1C3A5E}
\definecolor{pblue}{HTML}{3B78B5}
\definecolor{amber}{HTML}{F2B33D}
\definecolor{deepamber}{HTML}{D9902A}
\definecolor{pale}{HTML}{EEF3F9}
\definecolor{ink}{HTML}{2A2F36}
\definecolor{logonavy}{HTML}{173D63}
\definecolor{logodark}{HTML}{C98A1E}
\definecolor{logoamber}{HTML}{F2B446}
\definecolor{logocream}{HTML}{FEEFD2}
\definecolor{soft}{HTML}{5E6773}
\definecolor{agree}{HTML}{2F7D4F}
\definecolor{differ}{HTML}{B23B3B}
\color{ink}

\setlength{\parindent}{0pt}
\setlength{\parskip}{4pt}
\setlist{nosep,leftmargin=*,itemsep=2pt}
\renewcommand{\arraystretch}{1.15}

\pagestyle{fancy}
\fancyhf{}
\renewcommand{\headrulewidth}{0pt}
\fancyfoot[L]{\footnotesize\color{soft}\textbf{Pollis}\enspace New Methods\enspace·\enspace\nmtheme}
\fancyfoot[R]{\footnotesize\color{soft}\thepage}
\newcommand{\nmtheme}{}
\newcommand{\nmdate}{11 October 2026}

% The Pollis logo, redrawn so it prints sharply at any size
\newcommand{\pollislogo}[1]{%
  \begin{tikzpicture}[x=#1/512,y=-#1/512,baseline=(current bounding box.center)]
    \fill[logonavy] (0,0) rectangle (512,512);
    \fill[logodark] (36,422) -- (104,256) -- (172,422) -- cycle;
    \fill[logoamber] (104,422) -- (203,180) -- (302,422) -- cycle;
    \fill[logocream] (203,422) -- (339,89) -- (476,422) -- cycle;
  \end{tikzpicture}}

% Top of the first page: logo, name; then the kicker, title and a one-line promise
\newcommand{\nmhead}[3]{%
  \renewcommand{\nmtheme}{#1}%
  \thispagestyle{fancy}%
  \begin{minipage}[c]{13mm}\pollislogo{11mm}\end{minipage}%
  \begin{minipage}[c]{\dimexpr\linewidth-13mm\relax}
    {\color{navy}\Large\bfseries Pollis}\enspace{\color{soft}New Methods}\hfill
    {\color{deepamber}\small\bfseries\MakeUppercase{#1}}
  \end{minipage}\par\vspace{7mm}
  {\color{navy}\fontsize{22}{26}\selectfont\bfseries #2}\par\vspace{4pt}
  {\color{soft}\large #3}\par\vspace{5mm}}

\newcommand{\nmsection}[1]{\par\vspace{5pt}{\color{navy}\large\bfseries #1}\par\vspace{1pt}}
\newcommand{\nmsub}[1]{\par\vspace{2pt}{\color{deepamber}\bfseries #1}\par}
\newcommand{\gt}{\,{\color{deepamber}\bfseries ›}\,}
\newcommand{\menupath}[1]{{\bfseries\color{navy}#1}}
\newcommand{\fn}[1]{\texttt{\small #1}}
% Citations: \nmcite{key}{Breiman (2001)} links to the entry \nmref{key}{...} in the References
\newcommand{\nmcite}[2]{\hyperlink{ref:#1}{\color{pblue}#2}}
\newcommand{\nmref}[2]{\raisebox{1.1em}[0pt][0pt]{\hypertarget{ref:#1}{}}#2\par}
\newcommand{\nmdoi}[1]{\href{https://doi.org/#1}{\color{pblue}doi:#1}}
% Where a notebook is in Pollis: its title in Browse Methods, and its file
\newcommand{\browse}[2]{\menupath{Compose \gt New Methods \gt Browse Methods}: \emph{#1} (\fn{#2.ipynb})}

% Key facts at the top: source, missing in Julia, notebook, result
\newcommand{\factsheet}[4]{%
  \begin{tcolorbox}[enhanced, colback=pale, colframe=pale, boxrule=0pt, sharp corners,
    left=7pt, right=7pt, top=5pt, bottom=5pt, fontupper=\small]
  \begin{tabularx}{\linewidth}{@{}>{\leavevmode\bfseries\color{navy}}p{27mm}>{\leavevmode}X@{}}
  Tested against & #1\\
  Missing in Julia & #2\\
  Notebook & #3\\
  Result & #4
  \end{tabularx}
  \end{tcolorbox}}

% Boxes
\newtcolorbox{pbox}[2][]{enhanced, colback=pale, colframe=amber, boxrule=0pt,
  leftrule=3pt, sharp corners, left=7pt, right=7pt, top=6pt, bottom=6pt,
  fontupper=\small\raggedright, before upper={{\color{navy}\normalsize\bfseries #2}\par\smallskip}, #1}
\newtcolorbox{nbox}[2][]{enhanced, colback=navy, colframe=navy, boxrule=0pt,
  sharp corners, left=8pt, right=8pt, top=7pt, bottom=7pt, coltext=white,
  fontupper=\small\raggedright, before upper={{\color{amber}\normalsize\bfseries #2}\par\smallskip}, #1}

% Value tables: published beside Julia
\newenvironment{valuetable}[1]{%
  \par\small\tabularx{\linewidth}{#1}\toprule}{\bottomrule\endtabularx\par}
\newcommand{\ok}{{\color{agree}\faCheck}}
\newcommand{\no}{{\color{differ}\faTimes}}

% End of the document: credit and licence
\newcommand{\nmcredit}{%
  \par\vfill
  \begin{minipage}[b]{0.78\linewidth}\footnotesize\color{soft}
  Methods written and tested with Claude Code (Anthropic) in Pollis, from the published papers and their data.\par
  \textcopyright\ 2026 Antonio Saragga Seabra\enspace·\enspace Licensed under \href{https://creativecommons.org/licenses/by-sa/4.0/}{CC BY-SA 4.0}\enspace·\enspace Draft, \nmdate
  \end{minipage}\hfill
  \begin{minipage}[b]{0.18\linewidth}\raggedleft\pollislogo{9mm}\par\vspace{2pt}{\footnotesize contact@pollis.org}\end{minipage}}


\begin{document}

\nmhead{Volatility}
  {Periodic GARCH: volatility with a calendar dummy}
  {Does a shock after a weekend feed into tomorrow's risk as much as any other? Bollerslev and Ghysels (1996), reproduced.}

\factsheet
  {\nmcite{bollerslev1996}{Bollerslev and Ghysels (1996)}, \emph{Journal of Business \& Economic Statistics}, Table 2 (p.\,146): the 1,974 daily
   Deutsche mark / British pound returns of 1984--1991, with the paper's non-trading-day dummy (PollisDatasets \fn{bg96}).}
  {a dummy in the variance equation of a GARCH model, and the periodic GARCH model. ARCHModels.jl fits GARCH(1,1) only.}
  {\browse{Periodic GARCH: volatility with a calendar dummy}{pgarch}}
  {32 of 42 printed values to every digit, most of the rest in the last digit. Plain GARCH agrees with ARCHModels.jl to
   about $10^{-13}$.}

\nmsection{The problem}
Volatility is not the same every day. After a weekend or a holiday, news has piled up while markets were shut; a
GARCH(1,1) treats that day like any other. \nmcite{bollerslev1996}{Bollerslev and Ghysels (1996)} let the model change with a calendar dummy $s(t)$: the
mean, the level the variance returns to, and how much a shock feeds into the next day's variance can all differ when
$s(t) = 1$. The idea carries over to any periodic calendar: days of the week, months, trading sessions.

\nmsection{The method}
The mean is $\mu + \mu_1 s(t)$ and the conditional variance (p.\,146, eq.\ 14 and the note to Table 2) is
\[ \sigma^2_t = \omega + \omega_1 s_t + (\alpha + \alpha_1 s_t)\,(\varepsilon^2_{t-1} - \omega - \omega_1 s_{t-1})
   + \beta\,(\sigma^2_{t-1} - \omega - \omega_1 s_{t-1}). \]
$\omega + \omega_1 s(t)$ is the level the variance returns to, $\omega_1$ the extra variance when $s(t) = 1$, and $\alpha_1$
how much less, or more, a shock on such a day feeds forward. With $\alpha_1 = 0$ it is GARCH(1,1) with a dummy; with
$\omega_1 = \alpha_1 = \mu_1 = 0$, plain GARCH(1,1). \fn{pgarch} maximises the Gaussian likelihood and reports quasi-maximum
likelihood (sandwich) standard errors; any parameter can be held at zero. \fn{diagnostics} gives every statistic of
Table 2: information criteria, skewness and kurtosis of the standardised residuals, the statistics of \nmcite{ljung1978}{Ljung and Box (1978)} on them and on
their squares at 20 lags, and the paper's three loss functions.

\nmsection{What the tests show}
\begin{valuetable}{@{}lXl@{}}
\textbf{Test} & \textbf{What is compared} & \textbf{Agreement}\\\midrule
ARCHModels.jl & plain GARCH(1,1): estimates, robust s.e., log likelihood & to about $10^{-13}$\\
Table 2, col.\,5 & GARCH(1,1) with the dummy: 19 values & 16 to every digit, 3 in the last\\
Table 2, col.\,6 & P-GARCH(1,1): coefficients & every coefficient to 3 decimals\\
Table 2, col.\,6 & P-GARCH(1,1): s.e., SIC & 4 in the last digit\\
Table 2, col.\,6 & skewness $b_3$, kurtosis $b_4$ & $-.43$, 5.63 vs $-.46$, 5.87 (open)\\
\end{valuetable}

The last-digit differences, side by side (Julia against the paper):
\begin{itemize}
  \item Column 5: s.e.\ of the dummy .023 vs .022; $b_3$ $-.43$ vs $-.44$; $b_4$ 5.88 vs 5.89.
  \item Column 6: s.e.\ of $\omega$ .171 vs .169, of $\omega_1$ .027 vs .026, of $\alpha$ .042 vs .043; SIC $-2212.8$ vs $-2212.9$.
    These are what numerical Hessians of two programs give.
\end{itemize}

\nmsection{The economic result, reproduced}
On ordinary days shocks persist almost indefinitely: $\alpha + \beta = 1.0006$, which rounds to 1.001 against the paper's
1.000. After a non-trading day they die out much faster: $\alpha + \alpha_1 + \beta = .889$, as published. A Monday's shock
feeds less into Tuesday's variance, so the periodic model's volatility jumps less after weekends than plain GARCH's;
the notebook's example plots both over 1989--1990.

\nmsection{Still open}
The skewness and kurtosis of column 6: $-.43$ and 5.63 here against $-.46$ and 5.87 in print, while the Q statistics, MSE,
HMSE and log likelihood of the same standardised residuals all agree. The paper may compute $b_3$ and $b_4$ on a different
sample or residual definition. Not explained.

\nmsection{Further considerations}
\begin{itemize}
  \item \textbf{Several seasons.} The paper's general P-GARCH lets every parameter cycle with the calendar (five weekdays,
    twelve months). The notebook's single dummy is the case Table 2 tests; the general case needs one more index and no new
    ideas, and should come with a test of its own before it is offered.
  \item \textbf{Upstream.} A regressor in the variance equation is the missing piece in ARCHModels.jl; the dummy model
    is that plus a constraint. Column 5 is a ready-made test for such a contribution.
  \item \textbf{Asymmetry.} Calendar effects and leverage (EGARCH, GJR) are often studied together; a periodic GJR model
    would need a published table with its data before it becomes a new method here.
\end{itemize}

\nmsection{References}
{\footnotesize
\nmref{bollerslev1996}{Bollerslev, T. and Ghysels, E. (1996). Periodic autoregressive conditional heteroscedasticity. \emph{Journal of Business \& Economic Statistics} 14(2), 139--151. \nmdoi{10.1080/07350015.1996.10524640}}
\nmref{ljung1978}{Ljung, G.\,M. and Box, G.\,E.\,P. (1978). On a measure of lack of fit in time series models. \emph{Biometrika} 65(2), 297--303.}
\nmref{data}{Data: the returns as distributed with ARCHModels.jl, and the non-trading-day dummy as distributed with the R package rugarch (\fn{dmbp}).}
}

\nmcredit
\end{document}
